\documentclass[11pt,openany,oneside]{book}
\usepackage[T1]{fontenc}
\usepackage[margin=1in]{geometry}
\usepackage{hyperref}
\hypersetup{colorlinks=true,linkcolor=blue,urlcolor=blue}
\title{A Short History of the River Ferry}
\author{FlashTeX Fixtures Team}
\date{}
\begin{document}
\frontmatter
\tableofcontents
\phantomsection
\addcontentsline{toc}{chapter}{Preface}
{\Large\bfseries Preface\par}
The ferry has crossed the river since 1841, and this short book tells how. The story proper begins in Part~\ref{part:main} on page~\pageref{part:main}. These front pages carry roman numerals while the chapters ahead use arabic, and the table of contents above records both styles side by side once it is populated on the second pass.
\mainmatter
\part{The working ferry}
\label{part:main}
\chapter{From flatboat to steam}
\label{ch:boat}
The first crossing, previewed in the Preface, took eleven minutes against the current. The flatboat carried six passengers, two crates of tools, and the ferryman's dog, and within a year the run was hourly. The steam ferry replaced the flatboat in 1898: the crossing took four minutes and ran day and night, traffic peaked in 1911, and the logbooks record fog delays, ice delays, and one wedding held mid-river.
\end{document}
